\section{The nudge tax: user labor hidden in the chat}

When a model nearly completes a task, the user often supplies a series of small corrections. They restate the objective, name an overlooked constraint, provide a missing document, remind the product what it did earlier, open a website, copy a result into another app, or verify the output manually. Each intervention may feel trivial. In aggregate, the user is doing planning, integration, exception handling and quality assurance. The conversation transcript records the assistant's turns vividly but often fails to price the human work that made those turns useful. We call that unpriced work the \emph{nudge tax}.

The tax is not the mere existence of human involvement. Some approvals are necessary and valuable; a person may have to accept a consequential decision or confirm an external effect. The target is repeated tactical input that the system could have managed consistently once the task and authority were known. Figure~\ref{fig:nudge} gives an illustrative episode. Its spacing represents sequence, not time. The harness lane still includes a human approval; it removes routine reframing and state transfer only if evaluation shows that it can do so without increasing errors.

\begin{figure}[tbp]
\centering\begin{tikzpicture}[x=.92cm,y=1cm,>=Latex]
\draw[ink!25,line width=.7pt] (0,2.05)--(12,2.05);
\draw[ink!25,line width=.7pt] (0,.35)--(12,.35);
\node[anchor=east,font=\small\bfseries,text=ochre] at (-.15,2.05) {Manual nudging};
\node[anchor=east,font=\small\bfseries,text=teal] at (-.15,.35) {Harness};
\foreach \x in {0.5,2.7,4.9,7.1,9.3,11.5}{\fill[ochre] (\x,2.05) circle (2.5pt);}
\foreach \x in {0.5,5.0,9.5,11.5}{\fill[teal] (\x,.35) circle (2.5pt);}
\node[font=\scriptsize,align=center,text=ink] at (.5,2.56) {Request};
\node[font=\scriptsize,align=center,text=ink] at (2.7,2.56) {Reframe};
\node[font=\scriptsize,align=center,text=ink] at (4.9,2.56) {Supply\\context};
\node[font=\scriptsize,align=center,text=ink] at (7.1,2.56) {Move state\\by hand};
\node[font=\scriptsize,align=center,text=ink] at (9.3,2.56) {Approve};
\node[font=\scriptsize,align=center,text=ink] at (11.5,2.56) {Verify};
\node[font=\scriptsize,align=center,text=ink] at (.5,-.20) {Request};
\node[font=\scriptsize,align=center,text=ink] at (5,-.20) {Harness resolves\\routine friction};
\node[font=\scriptsize,align=center,text=ink] at (9.5,-.20) {Approve};
\node[font=\scriptsize,align=center,text=ink] at (11.5,-.20) {Verified};
\node[font=\scriptsize,text=ochre] at (6.1,1.42) {User-active touches};
\node[font=\scriptsize,text=teal] at (6.3,.85) {Automated state and exception handling};
\end{tikzpicture}

\caption{Two schematic interaction traces for the same permitted task. Dots are events, not measured minutes. A successful harness should reduce user-active touches while retaining approvals and independent verification.}
\label{fig:nudge}
\end{figure}

This distinction changes the baseline. Comparing a harness to a stock product used only once may show a large gain, but an expert user could obtain the same result through carefully sequenced prompting and manual transfer. The commercially relevant comparison includes that expert intervention, priced at the user's active time. A second comparison is needed for causality: same model, same task, and equal tool access, but different harness components. These answer different questions. The first asks whether a customer should pay a specialist. The second asks which design choice produced the gain.

The hidden labor also complicates claims about productivity. A field study of generative AI assistance in customer support measured changes in worker performance in a specific assisted-work setting \citep{genaiwork}; it is not evidence that an autonomous harness will complete account workflows. The productivity J-curve literature emphasizes that complementary organizational investments can precede observable gains from a general-purpose technology \citep{jcurve}. Both sources suggest that task-level measurement should include worker effort and integration cost, but neither supplies a refusal rate or ROI for the aftermarket opportunity proposed here.

An episode log should therefore count at least three classes of human time. \emph{Active nudging} is hands-on instruction, repair and data transfer. \emph{Required governance} is approval, identity confirmation or review that remains even in an excellent product. \emph{Waiting} is elapsed latency during which the user is not actively working but the task is unfinished. These classes have different economics. A harness may reduce active time while extending waiting; a faster model may reduce waiting while requiring more correction; a deterministic connector may eliminate both for a stable API. The reported ``time saved'' must specify which category changed.

The tax is also a source of product information. Repeated nudges reveal which facts the product failed to retain, which tool transitions it could not manage, and which success condition it failed to verify. A research prototype should preserve these interventions as structured observations, then ask which can be safely made repeatable. The goal is not to eliminate the user from every loop. It is to stop demanding a skilled operator every time a legitimate, recurring task crosses the same predictable gap.
